\subsection*{(1) Show that $ES_{\alpha}$ is monotone and cash-invariant}

\textbf{Monotonicity:}
Let $L_1$ and $L_2$ be two loss random variables such that $L_1 \le L_2$ almost surely. 
By the definition of Value at Risk, for any probability level $u \in (0,1)$, $VaR_u(L_1) \le VaR_u(L_2)$.
Using the integral representation of Expected Shortfall (which is directly equivalent to the conditional expectation for continuous distributions):
\begin{equation*}
ES_{\alpha}(L_1) = \frac{1}{1-\alpha}\int_{\alpha}^1 VaR_u(L_1) du \le \frac{1}{1-\alpha}\int_{\alpha}^1 VaR_u(L_2) du = ES_{\alpha}(L_2)
\end{equation*}
Thus, $ES_{\alpha}$ is monotone.

\textbf{Cash-invariance (Translation invariance):}
Let $c \in \mathbb{R}$ be a deterministic constant (e.g., adding cash to a portfolio reduces the loss by $c$). We evaluate $ES_{\alpha}(L + c)$.
First, note that $VaR_{\alpha}(L + c) = VaR_{\alpha}(L) + c$.
Using the conditional expectation formula provided:
\begin{align*}
ES_{\alpha}(L + c) &= \mathbb{E}[L + c \mid L + c \ge VaR_{\alpha}(L + c)] \\
&= \mathbb{E}[L + c \mid L + c \ge VaR_{\alpha}(L) + c] \\
&= \mathbb{E}[L + c \mid L \ge VaR_{\alpha}(L)] \\
&= \mathbb{E}[L \mid L \ge VaR_{\alpha}(L)] + c \\
&= ES_{\alpha}(L) + c
\end{align*}
Thus, $ES_{\alpha}$ is cash-invariant.


\subsection*{(2) Show that $ES_{\alpha}$ is positively homogeneous}

Let $\lambda > 0$ be a constant. We want to show $ES_{\alpha}(\lambda L) = \lambda ES_{\alpha}(L)$.
Since $\lambda > 0$, the function $f(x) = \lambda x$ is strictly increasing. 
From the results of Exercise 2, we know that $VaR_{\alpha}(\lambda L) = \lambda VaR_{\alpha}(L)$.
Substitute this into the provided formula:
\begin{align*}
ES_{\alpha}(\lambda L) &= \mathbb{E}[\lambda L \mid \lambda L \ge VaR_{\alpha}(\lambda L)] \\
&= \mathbb{E}[\lambda L \mid \lambda L \ge \lambda VaR_{\alpha}(L)] \\
&= \mathbb{E}[\lambda L \mid L \ge VaR_{\alpha}(L)] \quad \text{(since } \lambda > 0 \text{)} \\
&= \lambda \mathbb{E}[L \mid L \ge VaR_{\alpha}(L)] \\
&= \lambda ES_{\alpha}(L)
\end{align*}
Thus, $ES_{\alpha}$ is positively homogeneous.


\subsection*{(3) Prove the subadditivity of $ES_{\alpha}$}

To prove $ES_{\alpha}(L_1 + L_2) \le ES_{\alpha}(L_1) + ES_{\alpha}(L_2)$, we first establish a useful lemma using the secondary formula given in the exercise:
\begin{equation*}
ES_{\alpha}(L) = \frac{1}{1-\alpha} \mathbb{E}[L \cdot 1_{\{L \ge VaR_{\alpha}(L)\}}]
\end{equation*}

\textbf{Lemma:} For any event $A$ such that $\mathbb{P}(A) = 1 - \alpha$, we have $\mathbb{E}[L \cdot 1_A] \le \mathbb{E}[L \cdot 1_{\{L \ge q\}}]$ where $q = VaR_{\alpha}(L)$.

\textit{Proof of Lemma:} Let us look at the difference:
\begin{equation*}
\mathbb{E}[L \cdot 1_{\{L \ge q\}}] - \mathbb{E}[L \cdot 1_A] = \mathbb{E}[L (1_{\{L \ge q\}} - 1_A)]
\end{equation*}
Consider the value of $L (1_{\{L \ge q\}} - 1_A)$ inside the expectation:
\begin{itemize}
    \item If $L \ge q$, then $1_{\{L \ge q\}} = 1$. The term $(1_{\{L \ge q\}} - 1_A) \ge 0$. Since $L \ge q$, we have $L(1_{\{L \ge q\}} - 1_A) \ge q(1_{\{L \ge q\}} - 1_A)$.
    \item If $L < q$, then $1_{\{L \ge q\}} = 0$. The term $(1_{\{L \ge q\}} - 1_A) \le 0$. Since $L < q$ and the term is negative or zero, we also have $L(1_{\{L \ge q\}} - 1_A) \ge q(1_{\{L \ge q\}} - 1_A)$.
\end{itemize}
In both cases, $L(1_{\{L \ge q\}} - 1_A) \ge q(1_{\{L \ge q\}} - 1_A)$ almost surely.
Taking the expectation on both sides:
\begin{align*}
\mathbb{E}[L(1_{\{L \ge q\}} - 1_A)] &\ge q \mathbb{E}[1_{\{L \ge q\}} - 1_A] \\
&= q (\mathbb{P}(L \ge q) - \mathbb{P}(A)) \\
&= q ((1-\alpha) - (1-\alpha)) = 0
\end{align*}
This proves the lemma. It demonstrates that $ES_{\alpha}(L)$ can be viewed as the supremum of $\frac{1}{1-\alpha}\mathbb{E}[L \cdot 1_A]$ over all possible events $A$ with probability $1-\alpha$.

\textbf{Proof of Subadditivity:}
Let $L_P = L_1 + L_2$ and define the tail event for the combined portfolio $A^* = \{L_P \ge VaR_{\alpha}(L_P)\}$.
By our assumption, $\mathbb{P}(A^*) = 1 - \alpha$.
Using the formula for the portfolio:
\begin{align*}
ES_{\alpha}(L_1 + L_2) &= \frac{1}{1-\alpha} \mathbb{E}[(L_1 + L_2) 1_{A^*}] \\
&= \frac{1}{1-\alpha} \mathbb{E}[L_1 1_{A^*}] + \frac{1}{1-\alpha} \mathbb{E}[L_2 1_{A^*}]
\end{align*}
Since $\mathbb{P}(A^*) = 1 - \alpha$, we can apply our Lemma to $L_1$ and $L_2$ individually:
\begin{equation*}
\mathbb{E}[L_1 1_{A^*}] \le \mathbb{E}[L_1 1_{\{L_1 \ge VaR_{\alpha}(L_1)\}}] \quad \text{and} \quad \mathbb{E}[L_2 1_{A^*}] \le \mathbb{E}[L_2 1_{\{L_2 \ge VaR_{\alpha}(L_2)\}}]
\end{equation*}
Substituting these bounds back into the equation:
\begin{align*}
ES_{\alpha}(L_1 + L_2) &\le \frac{1}{1-\alpha} \mathbb{E}[L_1 1_{\{L_1 \ge VaR_{\alpha}(L_1)\}}] + \frac{1}{1-\alpha} \mathbb{E}[L_2 1_{\{L_2 \ge VaR_{\alpha}(L_2)\}}] \\
&= ES_{\alpha}(L_1) + ES_{\alpha}(L_2)
\end{align*}
Thus, Expected Shortfall is subadditive. Combining this with monotonicity, cash-invariance, and positive homogeneity, we conclude that $ES_{\alpha}$ satisfies all axioms of a coherent risk measure.